\rhead{CY501: Advanced Cryptography}
\phantomsection 
\section*{Advanced Cryptography (CY501)}
\label{major:CS_5}

\noindent \small{\hyperref[main:scheme]{$\leftarrow$ Back to Scheme}} \vspace{0.5cm}

\noindent \textbf{Credits:} 3 (3-0-0)

\subsection*{Course Content}
\noindent\textbf{Unit 1} \\
\textit{Block Cipher Cryptanalysis:} Differential cryptanalysis principles and attack complexity, Linear cryptanalysis and correlation matrices, Integral cryptanalysis and multidimensional characteristics, Truncated differential attacks, Impossible differential cryptanalysis, Boomerang and rectangle attacks, Slide attacks and self-similarity properties, Key schedule cryptanalysis techniques.

\vspace{0.3cm}

\noindent\textbf{Unit 2} \\
\textit{Stream Ciphers and Sequence Generators:} Linear Feedback Shift Registers (LFSRs) and Berlekamp-Massey algorithm, Nonlinear combination generators and correlation immunity, Stream cipher design criteria (period, linear complexity, balance), RC4 analysis (key scheduling biases, WEP attacks), Grain and Trivium stream ciphers, Related-key and distinguishing attacks on stream ciphers.

\vspace{0.3cm}

\noindent\textbf{Unit 3} \\
\textit{Public-Key Cryptosystems and Attacks:} RSA cryptanalysis (factoring algorithms - trial division, Fermat, quadratic sieve, number field sieve), Rabin cryptosystem and square roots modulo composite, Paillier homomorphic encryption, Elliptic Curve Discrete Logarithm Problem (ECDLP), Index calculus attack on elliptic curves, Pairing-based cryptography (bilinear maps, identity-based encryption).

\vspace{0.3cm}

\noindent\textbf{Unit 4} \\
\textit{Hash Function Design and Cryptanalysis:} Merkle-Damgård construction and length extension attacks, Sponge construction (Keccak/SHA-3), Hash function security properties (collision, preimage, second preimage resistance), Differential paths in hash functions, Boomerang distinguishers for compression functions, Multicollision attacks and Joux's attack, Hash-based signatures (Lamport, XMSS).

\vspace{0.3cm}

\noindent\textbf{Unit 5} \\
\textit{Post-Quantum Cryptography and Advanced Topics:} Quantum computing threats (Shor's algorithm, Grover's algorithm), Lattice-based cryptography (Learning With Errors - LWE, Ring-LWE), Code-based cryptography (McEliece cryptosystem), Hash-based signatures (SPHINCS+, LMS), Isogeny-based cryptography, Multivariate polynomial cryptography, NIST PQC standardization process and migration strategies.

\newpage
\rhead{Curriculum Schema}
